% Regression: t2_tcprimal — renders g_1 g_2 g_3^{-1} g_4^{-1} = e.
% Formula: g_1 g_2 g_3^{-1} g_4^{-1} = e .
% RMP 2011.12127, appendix "The Toric Code and quantum double models",
% eq:app:tcode-rep-primal (ref_ fig_app_tc-primal.pdf).
%
%   Quantum double for a finite group G: spins |g>, g in G, on the EDGES of
%   the oriented square lattice; the ground state is the equal-weight
%   superposition of configurations obeying the Gauss law at each vertex,
%       g_1 g_2 g_3^{-1} g_4^{-1} = e .
%   PEPS tensor = one blocked plaquette (aligned diagonally w.r.t. the
%   lattice).  Its four PHYSICAL indices are the edge spins h_1..h_4
%   (legs at the four sides); its four VIRTUAL indices sit at the
%   original-lattice vertices (the plaquette corners) and take the values
%       h_2 h_1^{-1},  h_3 h_2^{-1},  h_4 h_1^{-1},  h_3 h_4^{-1}
%   (the nonzero configurations).  The lines INSIDE the shaded tensor
%   indicate the original lattice (drawn dashed in the paper).
%
% Declarative model: a 3x3 oblique sheet; the shaded tensor is the region
% over all cells; the four physical legs are the boundary legs of the four
% mid-side sites (corners and centre removed); the original lattice inside
% is the corner-to-corner pair of distinguished edges.
\documentclass[border=6pt]{standalone}
\usepackage{tenkz}
\begin{document}
\begin{tenkzlattice}[rows=3, cols=3, frame=oblique, site=none, boundary legs]
  % the blocked-plaquette tensor (shaded)
  \tnregion[slot=selected]{(1-3,1-3)}
  % no beads, no interior bonds: corners + centre removed, so the only
  % lattice ink left is the four mid-side legs h_1 (west), h_2 (depth-up),
  % h_3 (east), h_4 (depth-down)
  \tnsite[removed]{(1,1)} \tnsite[removed]{(1,3)}
  \tnsite[removed]{(3,1)} \tnsite[removed]{(3,3)}
  \tnsite[removed]{(2,2)}
  % original lattice inside the tensor: the corner-to-corner cross carrying
  % the virtual Gauss-law indices h_{i+1} h_i^{-1}
  \tnedge[distinguished]{(1,1)-(3,3)}
  \tnedge[distinguished]{(1,3)-(3,1)}
\end{tenkzlattice}
\end{document}
